% Section 18.14, from lectures/L18/appendix/b_exercises.md; the requirements (Redovisning and
% Bedömning) from lectures/L18/README.md.

\section{Uppgifter}
\label{c18:sec:uppgifter}

\begin{aside}
\textbf{Så arbetar du med uppgifterna.} Varje uppgift löses i två steg. Först räknar du
\textbf{för hand} och skriver ned mellanstegen; det görs \textbf{före lektionen}. Sedan skriver
du under lektionen \textbf{ett program i C} som utför samma beräkning, och jämför resultaten.
Del~1--3 är obligatoriska och bedöms U/G. Del~4 och~5 är frivillig fördjupning för den som hinner
och påverkar inte bedömningen.

\facitnotis{L18}
\end{aside}

\begin{aside}
\note[Obs] Handberäkningarna för del~1--3 ska vara gjorda före lektionen och tas med dit.
Lektionstiden räcker inte till både beräkning och programmering.
\end{aside}

Koden skrivs och körs i OnlineGDB, \url{https://www.onlinegdb.com/online_c_compiler}, som i
\secref{c18:sec:miljo}. Kontrollera att språket är inställt på \textbf{C} och tryck på
\textbf{Run}. Vid felmeddelandet \code{undefined reference to 'sqrt'} saknas matematikbiblioteket;
lägg då till \code{-lm} under \textbf{Extra Compiler Flags} i inställningarna (kugghjulet).
Skriv ut flyttal med tre decimaler, exempelvis \code{printf("R = \%.3f ohm\n", r);}.

\subsection*{Arbetssätt: en fil per uppgift}
Samtliga uppgifter samlas i \textbf{ett och samma OnlineGDB-projekt}, med en fil per uppgift:
\begin{console}
assignment1_1.c   assignment2_1.c   assignment3_1.c
assignment1_2.c   assignment2_2.c   assignment3_2.c
assignment1_3.c   assignment2_3.c   assignment3_3.c
                  assignment2_4.c   main.c
\end{console}
Nya filer skapas med knappen \textbf{New file} i verktygsfältet, eller med \code{Ctrl + M}.

OnlineGDB kompilerar och länkar samtliga filer i projektet tillsammans. Därför får endast
\textbf{en} fil innehålla \code{main()}. Varje uppgift skrivs i stället som en egen funktion,
med samma namn som filen:
\begin{ccode}
/* assignment1_1.c */
#include <stdio.h>

void assignment1_1(void)
{
    printf("--- Uppgift 1.1 ---\n");
    /* Er kod här. */
}
\end{ccode}
I \file{main.c} deklareras uppgifterna överst och anropas sedan i tur och ordning. Lägg till en
rad per uppgift allteftersom du blir klar -- en funktion som deklareras men inte finns ger ett
länkfel:
\begin{ccode}
/* main.c */
void assignment1_1(void);
void assignment1_2(void);

int main(void)
{
    assignment1_1();
    assignment1_2();
    return 0;
}
\end{ccode}
Vid redovisningen räcker det då att trycka på \textbf{Run} en gång: samtliga uppgifter körs och
skriver ut sina resultat i ordning.

\begin{aside}
\note[Obs] Projektet sparas inte automatiskt. Kopiera filerna till din egen dator innan du
stänger fliken.
\end{aside}

\subsection*{Redovisning och bedömning}
Uppgifterna redovisas för läraren under lektionen:
\begin{itemize}
  \item \textbf{Handberäkningarna} för respektive uppgift, gjorda före lektionen och medtagna till
    den, handskrivna eller digitala och med tydliga steg.
  \item \textbf{Koden}: OnlineGDB-projektet med en fil per uppgift, som ska vara körbart och
    kommenterat. Samtliga uppgifter ska köras med en enda tryckning på \textbf{Run}.
  \item En \textbf{jämförelse} mellan handberäknat och beräknat resultat, med motivering av
    eventuella avvikelser.
\end{itemize}
Laborationen bedöms U/G och påverkar inte kursens betygspoäng. Den ingår i kravet att samtliga
lärandemål ska vara uppnådda för godkänt betyg på kursen.

\begin{narrowtable}{@{}lc@{}}
  \thead{Krav} & \thead{Bedömning} \\
  \midrule
  Del 1--3 genomförda och korrekt redovisade & G \\
  Del 1--3 ej genomförda eller ej redovisade & U \\
\end{narrowtable}

Del~4 och~5 är frivillig fördjupning och påverkar inte bedömningen.

% ------------------------------------------------------------------------------------------
\uppgiftsdel{Del 1}{Grunder i \code{math.h}}

\uppgift{1.1}{Resistans och effekt}
Två resistorer, $R_1 = 220\,\Omega$ och $R_2 = 330\,\Omega$, parallellkopplas över spänningen
$U = 12\,\text{V}$.
\begin{parts}
  \item Beräkna parallellresistansen $R_{\text{p}}$.
  \item Beräkna strömmen $I$ och den utvecklade effekten $P = \dfrac{U^2}{R_{\text{p}}}$.
  \item Skriv ett program som beräknar $R_{\text{p}}$, $I$ och $P$. Använd \code{pow()} för
    kvadreringen.
  \item Utöka programmet så att det även beräknar den spänning som krävs för att effekten ska bli
    $2{,}0\,\text{W}$, enligt $U = \sqrt{P R_{\text{p}}}$.
\end{parts}

\uppgift{1.2}{Decibel}
En förstärkare har inspänningen $U_{\text{in}} = 25\,\text{mV}$ och utspänningen
$U_{\text{ut}} = 1{,}6\,\text{V}$.
\begin{parts}
  \item Beräkna förstärkningen i dB.
  \item Ett filter dämpar signalen med $-3\,\text{dB}$. Beräkna kvoten
    $U_{\text{ut}}/U_{\text{in}}$ för filtret.
  \item Skriv ett program som beräknar båda värdena med \code{log10()} respektive \code{pow()}.
\end{parts}

\uppgift{1.3}{Vinklar och trigonometri}
En växelspänning ges av $u(t) = 325\sin(2\pi \cdot 50t)\,\text{V}$.
\begin{parts}
  \item Beräkna $u(t)$ vid $t = 2{,}0\,\text{ms}$.
  \item Ange vinkeln $2\pi \cdot 50t$ i både radianer och grader vid denna tidpunkt.
  \item Skriv ett program som beräknar $u(t)$ för $t = 0$, $2{,}0\,\text{ms}$, $5{,}0\,\text{ms}$
    och $10{,}0\,\text{ms}$. Skriv ut tiden, vinkeln i grader och spänningen på varje rad.
  \item Bestäm vinkeln för vektorn $\vec{v} = (-3, 4)$ med \code{atan2()}, i radianer och grader.
    Förklara varför \code{atan(4.0 / -3.0)} ger ett annat svar.
\end{parts}

% ------------------------------------------------------------------------------------------
\uppgiftsdel{Del 2}{Derivata och integraler}

\uppgift{2.1}{Numerisk derivering}
Låt $f(x) = x^3 - 2x + 1$.
\begin{parts}
  \item Derivera $f(x)$ analytiskt och beräkna $f'(2)$.
  \item Skriv en funktion \code{derivative()} som beräknar derivatan numeriskt med
    centraldifferensen
    \[
      f'(x) \approx \frac{f(x + h) - f(x - h)}{2h}.
    \]
  \item Beräkna $f'(2)$ numeriskt med $h = 10^{-2}$, $h = 10^{-6}$ och $h = 10^{-12}$. Skriv ut
    det absoluta felet mot det analytiska värdet för varje $h$.
  \item Vilket $h$ ger minst fel? Förklara varför både för stora och för små värden på $h$ ger
    sämre resultat.
\end{parts}

\uppgift{2.2}{Spänning över en spole}
Strömmen genom en spole ges av $i(t) = 5t^2\,\text{A}$. Spolens induktans är
$L = 0{,}20\,\text{H}$. Spänningen över spolen ges av
\[
  u(t) = L \cdot i'(t).
\]
\begin{parts}
  \item Bestäm $i'(t)$ analytiskt och därmed ett uttryck för $u(t)$.
  \item Beräkna $u(3)$ för hand.
  \item Skriv ett program som beräknar $u(3)$ genom att derivera $i(t)$ numeriskt, och jämför
    med handberäkningen.
\end{parts}

\uppgift{2.3}{Numerisk integrering}
Låt $f(x) = x^2$.
\begin{parts}
  \item Beräkna $\int_0^3 x^2\,dx$ analytiskt.
  \item Skriv en funktion \code{integral()} som beräknar den bestämda integralen med
    trapetsmetoden:
    \[
      \int_a^b f(x)\,dx \approx h\left[\frac{f(a) + f(b)}{2} + \sum_{k=1}^{n-1} f(a + kh)\right],
      \qquad h = \frac{b - a}{n}.
    \]
  \item Beräkna integralen med $n = 10$, $n = 100$ och $n = 1000$. Skriv ut det absoluta felet
    för varje $n$.
  \item Hur förändras felet när $n$ tiodubblas?
\end{parts}

\uppgift{2.4}{Laddning genom en ledare}
Strömmen ges av $i(t) = 3t + 2\,\text{A}$.
\begin{parts}
  \item Beräkna laddningen $q = \int_0^4 i(t)\,dt$ analytiskt.
  \item Beräkna samma laddning numeriskt med trapetsmetoden och $n = 4$.
  \item Varför blir det numeriska resultatet exakt i detta fall, men inte i \uppref{2.3}?
\end{parts}

% ------------------------------------------------------------------------------------------
\uppgiftsdel{Del 3}{Komplexa tal och fasorer}

\uppgift{3.1}{Räkning med komplexa tal}
Låt $z_1 = 4 + 3j$ och $z_2 = 2 - 5j$.
\begin{parts}
  \item Beräkna $z_1 + z_2$, $z_1 \cdot z_2$ och $\dfrac{z_1}{z_2}$ för hand.
  \item Beräkna $\abs{z_1}$ och $\arg(z_1)$ i radianer och grader.
  \item Skriv ett program som utför samtliga beräkningar med \file{complex.h}. Skriv ut varje
    resultat på formen \code{a + bj} med hjälp av \code{creal()} och \code{cimag()}.
\end{parts}

\uppgift{3.2}{Polär och rektangulär form}
Ett komplext tal ges på polär form som $z = 10\angle\frac{\pi}{6}$.
\begin{parts}
  \item Omvandla $z$ till rektangulär form för hand.
  \item Skriv ett program som utför omvandlingen med \code{cexp()}, och omvandlar tillbaka med
    \code{cabs()} och \code{carg()}.
  \item Kontrollera att du får tillbaka det ursprungliga talet. Varför bör resultaten jämföras med
    \code{fabs(a - b) < 1e-9} i stället för med \code{==}?
\end{parts}

\uppgift{3.3}{Impedans i en RLC-krets}
En seriekrets består av $R = 47\,\Omega$, $L = 100\,\text{mH}$ och $C = 10\,\upmu\text{F}$.
Kretsen matas med $U = 230\,\text{V}$ vid $f = 50\,\text{Hz}$.
\begin{parts}
  \item Beräkna $\omega$, $X_L = \omega L$ och $X_C = \dfrac{1}{\omega C}$.
  \item Ange den totala impedansen $Z = R + j(X_L - X_C)$ på rektangulär form.
  \item Beräkna $\abs{Z}$ och $\arg(Z)$ i grader. Är kretsen induktiv eller kapacitiv?
  \item Beräkna strömmens belopp $\abs{I} = \dfrac{\abs{U}}{\abs{Z}}$.
  \item Skriv ett program som beräknar allt ovanstående med \file{complex.h}. Låt $f$ vara en
    variabel, och kör programmet även för $f = 500\,\text{Hz}$. Vad händer med kretsens karaktär?
\end{parts}

% ------------------------------------------------------------------------------------------
\uppgiftsdel{Del 4}{Fördjupning (frivillig)}

\uppgift{4.1}{Sampling av en sinussignal}
En signal ges av $u(t) = 5\sin\bigl(2\pi \cdot 50t + \frac{\pi}{4}\bigr)\,\text{V}$ och samplas
med $f_s = 800\,\text{Hz}$.
\begin{parts}
  \item Kontrollera att samplingsteoremet är uppfyllt.
  \item Beräkna antalet sampel per period, $N = \dfrac{f_s}{f}$.
  \item Beräkna sampelvärdena $u_0$, $u_1$ och $u_2$ för hand enligt
    \[
      u_k = \abs{U} \sin\!\left(\frac{2\pi k}{N} + \delta\right).
    \]
  \item Skriv ett program som beräknar och skriver ut samtliga $N$ sampelvärden för en period,
    som CSV med en rad per sampel:
\begin{console}
k,u
0,3.536
1,4.619
\end{console}
  \item Kopiera utskriften till ett kalkylprogram och rita upp $u_k$ som funktion av $k$.
    Kontrollera att kurvan är sinusformad och att värdena upprepar sig efter $N$ sampel.
\end{parts}

\uppgift{4.2}{Addition av fasorer}
Tre spänningar med samma frekvens ges av fasorerna
\[
  U_1 = 10\angle 0, \qquad U_2 = 6\angle\frac{\pi}{3}, \qquad U_3 = 4\angle{-\frac{\pi}{2}}.
\]
\begin{parts}
  \item Omvandla samtliga fasorer till rektangulär form för hand.
  \item Beräkna summan $U = U_1 + U_2 + U_3$ på rektangulär form.
  \item Ange summan på polär form, $\abs{U}\angle\delta$, med $\delta$ i både radianer och
    grader.
  \item Skriv ett program som utför additionen med \file{complex.h} och skriver ut summan på både
    rektangulär och polär form.
  \item Rita fasorerna och deras summa i det komplexa talplanet.
\end{parts}

\uppgift{4.3}{Den sammansatta signalen}
\begin{parts}
  \item Skriv ett program som samplar den sammansatta signalen
    \[
      u(t) = u_1(t) + u_2(t) + u_3(t),
    \]
    där varje delsignal ges av motsvarande fasor i \uppref{4.2} och samtliga har frekvensen
    $f = 50\,\text{Hz}$.
  \item Skriv ut sampelvärdena som CSV med $f_s = 1600\,\text{Hz}$.
  \item Kontrollera att det största sampelvärdet överensstämmer med det $\abs{U}$ som beräknades
    i \uppref{4.2}. Motivera eventuell avvikelse.
\end{parts}

% ------------------------------------------------------------------------------------------
\uppgiftsdel{Del 5}{Extrauppgifter (frivillig)}
Uppgifterna nedan är extra träning för den som vill fördjupa sig ytterligare. De påverkar inte
bedömningen. Precis som i del~1--3 görs handberäkningen först och programmet därefter.

\uppgift{5.1}{Avrundning och absolutbelopp}
Betrakta talen $-2{,}7$, $3{,}2$ och $5{,}5$.
\begin{parts}
  \item Ange för hand vad \code{floor()}, \code{ceil()} och \code{round()} ger för vart och ett
    av talen.
  \item Skriv ett program som skriver ut alla nio värdena i en tabell.
  \item Varför ger \code{round(-2.5)} värdet $-3$ och inte $-2$?
  \item Vad är skillnaden mellan \code{fabs()} och \code{abs()}, och vilken bör användas i den
    här kursen?
\end{parts}

\uppgift{5.2}{Toleranskontroll av resistorer}
En resistor är märkt $4{,}7\,\text{k}\Omega \pm 5\,\%$.
\begin{parts}
  \item Beräkna det tillåtna intervallet för hand.
  \item Skriv funktionen
\begin{ccode}
int within_tolerance(double nominal, double tol_percent, double measured);
\end{ccode}
    som returnerar $1$ om det uppmätta värdet ligger inom toleransen, annars $0$.
  \item Testa funktionen med de uppmätta värdena $4600\,\Omega$, $4950\,\Omega$ och
    $4465\,\Omega$.
  \item Varför bör jämförelsen skrivas som
\begin{ccode}
fabs(measured - nominal) <= nominal * tol_percent / 100.0
\end{ccode}
    i stället för att jämföra mot två separata gränser?
\end{parts}

\uppgift{5.3}{Serie- och parallellresistans}
\begin{parts}
  \item Beräkna för hand $R_{\text{s}}$ och $R_{\text{p}}$ för $R_1 = 100\,\Omega$ och
    $R_2 = 400\,\Omega$.
  \item Skriv funktionerna
\begin{ccode}
double series(double r1, double r2)
double parallel(double r1, double r2)
\end{ccode}
  \item Skriv ett program som skriver ut båda värdena.
  \item Utöka programmet med en funktion som beräknar parallellresistansen för $n$ motstånd i en
    array, via summan av $1/R$. Testa med $\{100, 400, 200\}\,\Omega$.
\end{parts}

\uppgift{5.4}{Tabell över en RC-urladdning}
Spänningen ges av $u(t) = 12e^{-t/\tau}\,\text{V}$, där $\tau = RC$ med $R = 10\,\text{k}\Omega$
och $C = 100\,\upmu\text{F}$.
\begin{parts}
  \item Beräkna $\tau$ för hand.
  \item Beräkna $u(0)$, $u(1)$ och $u(2)$ för hand.
  \item Skriv ett program som med \code{exp()} skriver ut $u(t)$ för $t = 0$ till $5\,\text{s}$
    i steg om $0{,}5\,\text{s}$.
  \item Efter hur många tidskonstanter är spänningen under $1\,\%$ av startvärdet? Låt programmet
    leta upp svaret.
\end{parts}

\uppgift{5.5}{Andragradsekvation med alla tre fallen}
\begin{parts}
  \item Skriv ett program som löser $ax^2 + bx + c = 0$ med hjälp av diskriminanten.
  \item Programmet ska hantera de tre fallen $D > 0$, $D = 0$ och $D < 0$ separat.
  \item Testa med koefficienterna $(1, -5, 6)$, $(1, -4, 4)$ och $(1, 2, 5)$. Räkna ut de
    förväntade svaren för hand först.
  \item Skriv i fallet $D < 0$ ut rötterna på formen \code{a + bj} med hjälp av \code{sqrt(-D)}.
\end{parts}

\uppgift{5.6}{Numerisk ekvationslösning med intervallhalvering}
Sök roten till $f(x) = e^{-x} - x$ i intervallet $[0, 1]$.
\begin{parts}
  \item Beräkna $f(0)$ och $f(1)$ för hand och visa att de har olika tecken.
  \item Skriv en funktion \code{bisect()} som upprepat halverar intervallet tills dess bredd
    understiger $10^{-9}$.
  \item Kör programmet. Vilket värde på roten fås?
  \item Varför krävs det att $f(a)$ och $f(b)$ har olika tecken för att metoden ska fungera?
\end{parts}

\uppgift{5.7}{Effektivvärde genom numerisk integrering}
För en sinusspänning $u(t) = \abs{U}\sin(\omega t)$ gäller
\[
  U_{\text{RMS}} = \sqrt{\frac{1}{T}\int_0^T u(t)^2\,dt}.
\]
\begin{parts}
  \item Ange det teoretiska värdet på $U_{\text{RMS}}$ då $\abs{U} = 325\,\text{V}$.
  \item Skriv ett program som beräknar integralen numeriskt med trapetsmetoden och $n = 1000$
    delintervall över en period.
  \item Jämför med det teoretiska värdet och skriv ut det relativa felet.
  \item Vad händer med felet om $n$ minskas till $10$? Förklara varför.
\end{parts}
